\documentclass[../../main.tex]{subfiles}
\begin{document}
\section{Covariance technology: small-time operator-norm control}
\label{sec:covariance-tech}

Both Eldan subroutes meter the danger of localization through the same cut-free covariance
functional.
Recall from \eqref{eq:interface-def} the covariance excess $X_t=(\lmax(A_t)-1)_+$ and the
interface functional $\Xi_T(\mu)=\int_0^T\E X_t\dd t$; its second-moment companion is
\begin{equation}\label{eq:Xi2-def}
  \Xi^{(2)}_T(\mu)=\int_0^T\E X_t^2\dd t .
\end{equation}
The bootstrap of Section~\ref{sec:bootstrap} consumes $\Xi_T$ and the covariance reduction of
Section~\ref{sec:product-stress} consumes $\Xi^{(2)}_T$; both rest on the small-time
operator-norm control imported here. We isolate it as an assumption and discharge it against
the cited covariance technology. The exponent $C_2=2$ follows from the published sup-time
estimate below; the sharper $C_2=1$ conclusion is conditional on the July 2026 version-1
preprint input.

\begin{assumption}[Known small-time operator-norm control; matched to
{\cite{KlartagLehec2022Polylog,Letwin2026QuadraticKLS}}]
\label{hyp:KI}
There are universal constants $c_0,C_1,C_2>0$ such that for every isotropic log-concave $\mu$ on
$\R^n$, $n\ge3$,
\[
  \E\,\norm{A_t}_\op\ \le\ C_1
  \qquad\text{for all }0\le t\le t_1(n):=c_0(\log n)^{-C_2}.
\]
\end{assumption}

\noindent
This is the shape of the estimate underlying the Chen bootstrap as presented in Klartag's
lectures. The older route through the then-known Cheeger bound supplied $C_2=2$. Letwin's
dimension-free quadratic Poincar\'e inequality (Theorem~\ref{thm:letwin-qcts}) instead bounds the
third-moment parameter $\kappa_n$ universally; inserted into the precise Klartag--Lehec moment
window, it gives $C_2=1$. The standard two-sided-exponential product example suggests that
$1/\log n$ is the natural endpoint for covariance-only control.

\begin{theorem}[{Klartag--Lehec; the sup-over-time form is \cite[Thm.~61]{KLnotes}}]
\label{thm:KL-window}
There is a universal constant $C$ such that for every isotropic log-concave $\mu$ on $\R^n$
and every $t\le\dfrac1{C\log^2n}$,
\[
  \Prob\bigl(\exists\,s\le t:\ \norm{A_s}_\op\ge2\bigr)\ \le\ \exp\Bigl(-\frac1{Ct}\Bigr).
\]
\end{theorem}

The newer parallel-coupling preprint contains rank-sensitive information that is stronger than
an operator-norm window but still cut-free.  We record it without wiring it to a live gate.

\begin{theorem}[{Stopped rank tails; \cite{KlartagLehec2025ThinShell}}]
\label{thm:kl-stopped-rank-tail}
For the simplified stochastic localization of an isotropic compactly supported log-concave law,
write $\lambda_1(t)\ge\cdots\ge\lambda_n(t)$ for the eigenvalues of $A_t$.  There is a universal
$C$ such that, for every stopping time $\sigma$ and $t>0$,
\[
  \sum_{i=1}^n\Prob\bigl(\lambda_i(t\wedge\sigma)\ge3\bigr)
  \le Cn\exp(-t^{-1/8}).
\]
Consequently, if $\sigma_k=\inf\{t:\lambda_k(t)\ge3\}$, then
\[
  \Prob(\sigma_k\le t)\le C\frac nk\exp(-t^{-1/8}),
  \qquad
  \E\sigma_k^{-2}\le C\left(1+\log\frac nk\right)^{16}.
\]
\end{theorem}

\begin{theorem}[{Integrated rank covariance; \cite{KlartagLehec2025ThinShell}}]
\label{thm:kl-integrated-rank-covariance}
Under the same hypotheses,
\[
  \E\sum_{i=1}^n
  \exp\left(2\int_0^1\lambda_i(t)\dd t\right)\le Cn .
\]
\end{theorem}

Both results are imported from version 2 of an unreviewed preprint.  They control how many
covariance eigenvalues are large and how long each rank can remain large.  They do not control
the orientation of a cut tensor $K_t$, a posterior Hessian $H_t$, or an eigenfunction source
relative to those eigenspaces.  In particular, neither theorem by itself discharges
Question~\ref{q:upgrade}, Question~\ref{q:mm-spectral-occupation}, or the high-rank part of
Question~\ref{q:stein-weighted}.

\begin{proposition}[Letwin's third-moment bound]
\label{prop:letwin-kappa}
Define
\[
  \kappa_n
  =\sup_{\nu}\sup_{\theta\in S^{n-1}}
    \norm{\E_{X\sim\nu}\bigl[\inner{X}{\theta}\,X\otimes X\bigr]}_{\HS},
\]
where the first supremum is over isotropic log-concave laws on $\R^n$. Conditional on the
preprint input of Theorem~\ref{thm:letwin-qcts},
\[
  \kappa_n\le2\sqrt2.
\]
\end{proposition}

\begin{proof}
Fix $\nu,\theta$ and put $M=\E[\inner{X}{\theta}\,X\otimes X]$. Isotropy and centering give
\[
  \norm M_{\HS}^2
  =\E\!\left[\inner{X}{\theta}\bigl(X^TMX-\Tr M\bigr)\right].
\]
Cauchy--Schwarz and Theorem~\ref{thm:letwin-qcts} therefore imply
$\norm M_{\HS}^4\le8\norm M_{\HS}^2$. The claim follows, including the case $M=0$.
\end{proof}

\begin{corollary}[The quadratic-chaos covariance window]
\label{cor:letwin-window}
For every fixed $p\ge1$ there are universal constants $c,C_p>0$ such that, for every isotropic
log-concave initial measure and
\[
  0\le t\le \frac{c}{\log n},
  \qquad
  \E\norm{A_t}_{\op}^p\le C_p.
\]
\end{corollary}

\begin{proof}
Klartag--Lehec \cite[Cor.~5.4]{KlartagLehec2022Polylog} prove the displayed moment bound on
$t\le(C\kappa_n^2\log n)^{-1}$. Apply Proposition~\ref{prop:letwin-kappa}.
\end{proof}

\begin{remark}[Attribution and the 2026 window upgrade]
\label{rem:kl-window-verified}
Theorem~\ref{thm:KL-window} --- \emph{sup over time}, threshold $2$, clean window
$t\le(C\log^2n)^{-1}$ with no thin-shell constant --- is \cite[Thm.~61]{KLnotes}. It should
\emph{not} be attributed to \cite[Lemma~5.2]{KlartagLehec2022Polylog}, which is the weaker
\emph{fixed-time} bound $\Prob(\norm{A_t}_\op\ge2)\le e^{-1/(Ct)}$ on the
$\kappa_n$-dependent window $t\le(C\kappa_n^2\log n)^{-1}$. The second moment is moreover
published in exactly the strength needed:
\cite[Cor.~5.4]{KlartagLehec2022Polylog} gives $\E\norm{A_t}_\op^p\le C_p$ for every $p\ge1$
on that window, so $\E X_t^2\le\E\norm{A_t}_\op^2\le C$ directly. Proposition~\ref{prop:letwin-kappa}
now makes the latter window $c/\log n$, resolving the former intermediate-window question at the
fixed-time moment level. These statement numbers have
been verified against arXiv:2406.01324v2 (the version matching the published
Bull.\ Amer.\ Math.\ Soc.\ \textbf{62} (2025), no.~4, article) for \cite{KLnotes}, and against
arXiv:2203.15551 for \cite{KlartagLehec2022Polylog}.
\end{remark}

\begin{corollary}[Discharge of Assumption~\ref{hyp:KI}]
\label{cor:KI-discharged}
The published Theorem~\ref{thm:KL-window}, together with the Brascamp--Lieb cap, discharges
Assumption~\ref{hyp:KI} with $C_2=2$.
\end{corollary}

\begin{proof}
For $t\le c/\log^2n$, split according to $\norm{A_t}_\op<2$ and use
Theorem~\ref{thm:KL-window} plus $\norm{A_t}_\op\le t^{-1}$ to obtain
$\E\norm{A_t}_\op\le2+t^{-1}e^{-1/(Ct)}\le C_1$.
\end{proof}

\begin{corollary}[Letwin-v1 sharpening of the covariance window]
\label{cor:KI-letwin}
Conditional on the preprint input of Theorem~\ref{thm:letwin-qcts},
Assumption~\ref{hyp:KI} holds with $C_2=1$: there is a universal constant $c_0$ such
that for every isotropic log-concave $\mu$ and every $t\le c_0(\log n)^{-1}$,
\[
  \E\,\norm{A_t}_\op\ \le\ C_1 .
\]
Consequently, conditional on the same version-1 preprint input, the polylog evaluation of the
interface functional ($\Xi_T\le C(1+\log\log n)$, Corollary~\ref{cor:loglog}) may use the
larger cutoff $t_1(n)=c_0/\log n$.  This improves the verified early-time window but does not
remove the $\log\log n$ universal-time loss or any of the geometric inputs downstream.
\end{corollary}

\begin{proof}
This is Corollary~\ref{cor:letwin-window} with $p=1$.
\end{proof}
\end{document}
